\section{在线参数辨识}
第3节的闭环制导律\eqref{eq:48}及其积分增益\eqref{eq:34}--\eqref{eq:37}依赖投影系数${b_i},{b_T}$、视线角速度合成量${\dot Q_i},{\dot Q_T}$以及NCV机动在视线上的投影${v_{ri}}$。在理想信息下，这些量可由式\eqref{eq:13}与\eqref{eq:15}解析计算；但在\cref{asm:meas-noise}的量测条件下，将含噪的视线角与航迹角直接代入解析式，会使制导指令产生有偏且相互相关的误差。为此，本节不再依赖解析式，而是由相对距离运动学在线辨识状态空间模型\eqref{eq:12}与\eqref{eq:14}中的非平凡元素，为制导律提供其所需参数的估计。

\subsection{参数辨识问题构建}
根据状态空间方程\eqref{eq:12}和\eqref{eq:14}，每个方程的第二行包含非平凡元素。对于规避者–NCV对，从方程\eqref{eq:10}中提取二阶动力学：

\begin{equation}
\label{eq:49}
{\ddot r_{ADi}}(t) = \underbrace {(\dot q_{yi}^2 + \dot q_{zi}^2{{\cos }^2}{q_{yi}})}_{A_{21}^{(i)}} \cdot {r_{ADi}}(t) + {B_{1,i}} \cdot {a_{yA}}(t) + {B_{2,i}} \cdot {a_{zA}}(t) + {C_{1i}}(t)
\end{equation}
同理，对规避者–目标对，由式\eqref{eq:11}可得

\begin{equation}
\label{eq:50}
{\ddot r_{AT}}(t) = \underbrace {(\dot q_\theta ^2 + \dot q_\psi ^2{{\cos }^2}{q_\theta })}_{A_{21}^{(T)}} \cdot {r_{AT}}(t) + {B_{1,T}} \cdot {a_{yA}}(t) + {B_{2,T}} \cdot {a_{zA}}(t)
\end{equation}
式\eqref{eq:49}–\eqref{eq:50}对未知参数均为线性。对比式\eqref{eq:49}和式\eqref{eq:3}可知${C_{1i}} =  - {v_{ri}}$，即第i枚NCV加速度在与规避者视线上投影的负值。

所有信号在制导周期${t_k} = k\Delta t$内采样。基于${t_k}$时刻参数估计计算的加速度指令在区间$[{t_k},{t_{k + 1}})$内保持不变。由于在${t_k}$时刻重构的输出应与前一采样区间内实际施加的加速度相配对，而不应与${t_k}$时刻正在计算的指令相配对，故定义规避者–NCV对的因果回归（causal regression）模型：

\begin{equation}
\label{eq:51}
{y_i}(k) = \phi _i^T(k - 1){\kern 1pt} {\vartheta _i} + {e_i}(k)
\end{equation}
其中

\begin{equation}
\label{eq:52}
{y_i}(k) = \hatddot{r}_{ADi}(k),\quad {\vartheta _i} = \left[ {\begin{array}{*{20}{c}}
  {A_{21}^{(i)}} \\ 
  {{B_{1,i}}} \\ 
  {{B_{2,i}}} \\ 
  {{C_{1i}}} 
\end{array}} \right],\quad {\phi _i}(k - 1) = \left[ {\begin{array}{*{20}{c}}
  {{r_{ADi}}(k - 1)} \\ 
  {{a_{yA}}(k - 1)} \\ 
  {{a_{zA}}(k - 1)} \\ 
  1 
\end{array}} \right]
\end{equation}
对规避者–目标对：

\begin{equation}
\label{eq:53}
{y_T}(k) = \phi _T^T(k - 1){\kern 1pt} {\vartheta _T} + {e_T}(k)
\end{equation}
其中

\begin{equation}
\label{eq:54}
{y_T}(k) = \hatddot{r}_{AT}(k),\quad {\vartheta _T} = \left[ {\begin{array}{*{20}{c}}
  {A_{21}^{(T)}} \\ 
  {{B_{1,T}}} \\ 
  {{B_{2,T}}} 
\end{array}} \right],\quad {\phi _T}(k - 1) = \left[ {\begin{array}{*{20}{c}}
  {{r_{AT}}(k - 1)} \\ 
  {{a_{yA}}(k - 1)} \\ 
  {{a_{zA}}(k - 1)} 
\end{array}} \right]
\end{equation}
式\eqref{eq:52}和\eqref{eq:54}中，${a_{yA}}(k - 1)$和${a_{zA}}(k - 1)$为在区间$[{t_{k - 1}},{t_k})$内实际施加的限幅后加速度，$\hatddot{r}_{ADi}(k)$、$\hatddot{r}_{AT}(k)$分别为采用卡尔曼滤波得到的规避者与第i枚NCV、目标的预估相对加速度。式\eqref{eq:51}和\eqref{eq:53}中的${\vartheta _i}$、${\vartheta _T}$为等效参数，其真值随视线几何连续变化：单条回归方程在区间$[{t_{k - 1}},{t_k})$内将其视为常数，式\eqref{eq:55}进而以同一常向量拟合遗忘窗内的全部数据，参数的慢变由遗忘因子$\rho $予以跟踪；在${t_k}$时刻递推得到的估计$\hat{\vartheta}_{i}(k)$、$\hat{\vartheta}_{T}(k)$即为当前可用的最新估计。残差$e(k)$表示回归残差，主要包括采用Kalman滤波根据距离和距离率估计相对加速度所引入的误差，以及以单一常向量表示遗忘窗内时变参数所产生的建模误差。

结构化辨识利用状态矩阵已知的稀疏模式，即${{\mathbf{A}}_i},{{\mathbf{A}}_T}$的$(1,1),(1,2){\text{,}}(2,2)$元素分别固定为0, 1, 0，${{\mathbf{B}}_i},{{\mathbf{B}}_T}$的第一行均为零。仅估计非平凡元素，从而降低参数维数并保证物理一致性。待辨识的非平凡元素汇总于表~\ref{tab:nontrivial}。

\begin{table}[htbp]
\centering
\caption{待辨识非平凡元素}
\label{tab:nontrivial}
\small
\begin{tabularx}{\linewidth}{lXXX}
\toprule
方程 & 矩阵 & 元素 & 物理含义 \\
\midrule
式\eqref{eq:12} & ${{\mathbf{A}}_i}$ & $A_{21}^{(i)} = \dot q_{yi}^2 + \dot q_{zi}^2{\cos ^2}{q_{yi}}$ & 视线角速率耦合 \\
式\eqref{eq:12} & ${{\mathbf{B}}_i}$ & ${B_{1,i}},{B_{2,i}}$ & 规避者加速度投影系数（NCV方向） \\
式\eqref{eq:12} & ${{\mathbf{C}}_i}$ & ${C_{1i}}$ & NCV机动加速度偏置 \\
式\eqref{eq:14} & ${{\mathbf{A}}_T}$ & $A_{21}^{(T)} = \dot q_\theta ^2 + \dot q_\psi ^2{\cos ^2}{q_\theta }$ & 目标-规避者视线角速率耦合 \\
式\eqref{eq:14} & ${{\mathbf{B}}_T}$ & ${B_{1,T}},{B_{2,T}}$ & 规避者加速度投影系数（目标方向） \\
\bottomrule
\end{tabularx}
\end{table}
\FloatBarrier

在\cref{asm:meas-noise}的量测条件下，$\ddot r$须由$r$与$\dot r$在线估计，回归残差$e$非零，且部分自由参数（尤其是投影系数${B_{1,i}},{B_{2,i}},{C_{1i}}$以及${B_{1,T}},{B_{2,T}}$）在某些几何中可能被弱激发。为提高弱激励条件下参数辨识的稳定性，并抑制相对加速度估计误差引起的伪参数分量，在带遗忘因子的最小二乘代价函数中附加${\ell _1}$(Lasso)正则项。对于规避者–NCV对$(i = 1,2)$

\begin{equation}
\label{eq:55}
{J_i}({\vartheta _i}) = \sum\limits_{k = 1}^N {{\rho ^{N - k}}} {\left[ {{y_i}(k) - \phi _i^T(k - 1){\vartheta _i}} \right]^2} + \mu {\left\| {{\vartheta _i}} \right\|_1}
\end{equation}
对规避者-目标对同样构造${J_T}({\vartheta _T})$，其中$\rho  \in (0,1]$为遗忘因子，$\mu  \geqslant 0$为稀疏度权重。由于$\phi $的各分量具有不同的物理量纲，实际实现中采用对角缩放矩阵$M$作归一化：

\begin{equation}
\label{eq:56}
\bar{\phi} = {M^{ - 1}}\phi ,\quad {\vartheta ^{{\text{sc}}}} = M\vartheta ,\quad {M_i} = {\text{diag}}({r_0},{a_0},{a_0},1),\quad {M_T} = {\text{diag}}({r_0},{a_0},{a_0})
\end{equation}
其中，${r_0}$和${a_0}$分别为标称距离尺度和加速度尺度。该变换保持预测量不变，即$\bar{\phi}^{T}{\vartheta ^{{\text{sc}}}}{\text{ = }}{\phi ^T}\vartheta $。在以下递推中，为简洁起见省略上标“sc”，所有变量均默认为归一化变量。

\begin{problem}[结构稀疏参数辨识]
\label{prb:id}
给定\cref{asm:meas-noise}所述导引头相对距离、距离率测量与规避者实际施加的加速度，在已知状态矩阵稀疏结构的约束下，以式\eqref{eq:55}的带遗忘因子${\ell _1}$正则准则为目标，在每个制导周期$\Delta t$内仅利用${t_k}$及以前的量，递推估计${\vartheta _i}$（$i = 1,2$）与${\vartheta _T}$，为制导律提供${{\mathbf{A}}_i},{{\mathbf{B}}_i},{{\mathbf{C}}_i},{{\mathbf{A}}_T},{{\mathbf{B}}_T}$的在线估计。
\end{problem}
\subsection{带${\ell _1}$近端更新的稀疏FF-RLS}
为求解\cref{prb:id}，本小节构造带${\ell _1}$近端更新的稀疏FF-RLS估计器：在每个制导周期内，先由遗忘因子递推最小二乘完成二次更新，再以一次逐元素软阈值实现式\eqref{eq:55}中${\ell _1}$正则所对应的稀疏收缩，并将${A_{21}}$的估计投影至物理可行域，从而改善弱激励与含噪残差下时变参数的跟踪性能。遗忘因子$\rho$对旧数据作指数折扣，使估计器能适应交会过程中视线几何关系的持续变化。对规避者–NCV对$(i = 1,2)$，其二次更新为：

\begin{equation}
\label{eq:57}
{{\mathbf{K}}_i}(k) = \frac{{{{\mathbf{P}}_i}(k - 1){\kern 1pt} {\phi _i}(k - 1)}}{{\rho  + \phi _i^T(k - 1){\kern 1pt} {{\mathbf{P}}_i}(k - 1){\kern 1pt} {\phi _i}(k - 1)}}
\end{equation}
\begin{equation}
\label{eq:58}
\bar{\vartheta}_{i}(k) = \hat{\vartheta}_{i}(k - 1) + {{\mathbf{K}}_i}(k)[{y_i}(k) - \phi _i^T(k - 1){\kern 1pt} \hat{\vartheta}_{i}(k - 1)]
\end{equation}
\begin{equation}
\label{eq:59}
{{\mathbf{P}}_i}(k) = \frac{1}{\rho }[{\mathbf{I}} - {{\mathbf{K}}_i}(k){\kern 1pt} \phi _i^T(k - 1)]{\kern 1pt} {{\mathbf{P}}_i}(k - 1)
\end{equation}
其中${{\mathbf{K}}_i}(k)$为增益矩阵，${{\mathbf{P}}_i}(k)$为估计误差协方差矩阵，$\bar{\vartheta}_{i}(k)$为二次更新得到的中间估计。规避者–目标对$\bar{\vartheta}_{T}(k)$的更新方程形式相同，只需将各项下标i替换为T。在RLS更新后，通过一次逐元素软阈值操作实现${\ell _1}$正则对应的稀疏收缩：

\begin{equation}
\label{eq:60}
\hat{\vartheta}_{i}(k) = {\mathcal{S}_\eta }(\bar{\vartheta}_{i}(k))
\end{equation}
其中，${[{\mathcal{S}_\eta }(v)]_j} = {\text{sign}}({v_j})\max (|{v_j}| - \eta ,0)$，$\eta  \geqslant 0$为归一化参数域中的软阈值，与式\eqref{eq:55}中的稀疏度权重$\mu $相对应，其取值由离线仿真标定。

对$\bar{\vartheta}_{T}(k)$采用相同操作。$\hat{\vartheta}_{i}(k)$为第$k$步最终估计，下一周期式\eqref{eq:58}中的$\hat{\vartheta}_{i}(k - 1)$取该阈值后的值。为兼顾在线计算效率，该单步收缩用于抑制含噪与弱激励条件下出现的无效小幅参数分量，并不等价于在每个采样时刻精确求解式\eqref{eq:55}。由于${A_{21}}$是平方项之和，并通过$\dot Q = \sqrt {{A_{21}}} $进入制导律，其估计值应投影到物理可行域：

\begin{equation}
\label{eq:61}
\hat A_{21}^{(c)}(k) \leftarrow {\text{max}}(\hat A_{21}^{(c)}(k),0),\quad c \in \{ 1,2,T\}
\end{equation}
参数估计值由式\eqref{eq:13}和式\eqref{eq:15}的解析表达式初始化，初始值由初始测量值计算并经$M$缩放后进入递推；同时令${{\mathbf{P}}_i}(0) = {\delta ^{ - 1}}{\mathbf{I}}$和${{\mathbf{P}}_T}(0) = {\delta ^{ - 1}}{\mathbf{I}}$，其中$\delta  \ll 1$。对于NCV通道，${C_{1i}}$的初始值求解需采用由PN模型式\eqref{eq:9}得到的NCV加速度。

按\cref{asm:meas-noise}的量测条件，相对加速度$\ddot r$无法直接量测。由于对离散距离率直接进行数值差分会引入离散化误差，因此对每个距离通道采用三状态常加速度Kalman滤波器：

\begin{equation}
\label{eq:62}
\xi (k) = {[r(k),\dot r(k),\ddot r(k)]^T},\quad {z_\xi }(k) = {[r(k),\dot r(k)]^T}
\end{equation}
滤波器转移状态矩阵为\[{{\mathbf{F}}_\xi } = \left[ {\begin{array}{*{20}{c}}
  1&{\Delta t}&{\tfrac{1}{2}\Delta {t^2}} \\ 
  0&1&{\Delta t} \\ 
  0&0&1 
\end{array}} \right]\]，量测矩阵为\[{{\mathbf{H}}_\xi } = \begin{bmatrix}
  1 & 0 & 0 \\
  0 & 1 & 0
\end{bmatrix}\]，并将相对加加速度建模为白噪声。滤波器后验估计的第三个分量作为回归输出，即$y(k) = [0,0,1]\hat{\xi}(k|k)$

最后，将辨识得到的参数重新组装为状态空间矩阵：

\begin{equation}
\label{eq:63}
\hat{\mathbf{A}}_{i} = \left[ {\begin{array}{*{20}{c}}
  0&1 \\ 
  {\hat A_{21}^{(i)}}&0 
\end{array}} \right],\quad \hat{\mathbf{B}}_{i} = \left[ {\begin{array}{*{20}{c}}
  0&0 \\ 
  {{{\hat B}_{1,i}}}&{{{\hat B}_{2,i}}} 
\end{array}} \right],\quad \hat{\mathbf{C}}_{i} = \left[ {\begin{array}{*{20}{c}}
  0 \\ 
  {{{\hat C}_{1i}}} 
\end{array}} \right]
\end{equation}
\begin{equation}
\label{eq:64}
\hat{\mathbf{A}}_{T} = \left[ {\begin{array}{*{20}{c}}
  0&1 \\ 
  {\hat A_{21}^{(T)}}&0 
\end{array}} \right],\quad \hat{\mathbf{B}}_{T} = \left[ {\begin{array}{*{20}{c}}
  0&0 \\ 
  {{{\hat B}_{1,T}}}&{{{\hat B}_{2,T}}} 
\end{array}} \right]
\end{equation}
式\eqref{eq:63}--\eqref{eq:66}中各估计量均已按$\vartheta  = {M^{ - 1}}{\vartheta ^{{\text{sc}}}}$反变换回物理量纲。制导律关于第i枚NCV所需的量由下式获得：

\begin{equation}
\label{eq:65}
\hat{\dot Q}_{i} = \sqrt {\hat A_{21}^{(i)}} ,\quad \hat{b}_{i} = {\left[ {{{\hat B}_{1,i}},{{\hat B}_{2,i}}} \right]^T},\quad \hat{v}_{ri} =  - \hat{C}_{1i}
\end{equation}
关于目标的所需的量由下式获得：

\begin{equation}
\label{eq:66}
\hat{\dot Q}_{T} = \sqrt {\hat A_{21}^{(T)}} ,\quad \hat{b}_{T} = {\left[ {{{\hat B}_{1,T}},{{\hat B}_{2,T}}} \right]^T}
\end{equation}
这些估计量替代解析计算得到的对应量，用于最优制导律\eqref{eq:48}与积分增益\eqref{eq:34}-\eqref{eq:37}计算，形成闭环辨识–制导架构。
\begin{remark}
上述辨识器的精度受以下因素影响。遗忘窗较短时，${r}$在窗内变化平缓，与常数回归项不易区分；而${a_{yA}}$、${a_{zA}}$由闭环制导律产生，因此各参数能否被充分辨识取决于交会过程中机动的丰富程度。回归输出$\hatddot{r}$由常加速度Kalman滤波估计，其滞后会在参数估计中引入一定偏差。此外，末段${C_{1i}}$随NCV比例导引加速度快速变化，将其在一个制导周期内视为常数的近似随之变弱。
\end{remark}

